\documentclass[10pt,a4paper]{amsart}
\usepackage[margin=19mm]{geometry}
\usepackage{amsmath,amssymb,mathtools}
\usepackage[scr=rsfs,cal=euler]{mathalfa}
\usepackage{xfrac}
\usepackage[unicode,hidelinks,bookmarks=false]{hyperref}
\newcommand{\ie}{i.e.\ }
\newcommand{\cf}{cf.\ }
\newcommand{\resp}{respectively\ }
\newcommand{\Subsection}{\subsection}
\providecommand{\nobreakdash}{\nobreak\hbox{-}}
\setlength{\emergencystretch}{2em}
\multlinegap0pt
\usepackage{amssymb}
\newtheorem{lemma}{Lemma}
\newtheorem{theorem}{Theorem}
\title{A Structure Sheaf for Kirch Topology}
\author{Alexander Borisov}
\date{Source: arXiv:2605.05515v1 (CC BY 4.0)}
\begin{document}
\maketitle

\noindent\emph{Source and licence.} A Structure Sheaf for Kirch Topology. Version arXiv:2605.05515v1 (CC BY 4.0), distributed by arXiv under the Creative Commons Attribution 4.0 International licence, http://creativecommons.org/licenses/by/4.0/, as recorded in the arXiv licence metadata for this exact version and shown on the abstract page https://arxiv.org/abs/2605.05515v1. Full licence notice: \texttt{licenses/arxiv-2605.05515-CC-BY-4.0.txt}.\par\smallskip

\noindent\textit{Editorial scope.} Counted unit: the article's theorem MainH1 (source0.tex lines 395--396). The article's complete original proof of it is delivered with it (source0.tex lines 398--417, one \texttt{proof} environment of 20 lines). No mathematics is written, repaired or re-derived: the statement and the proof are the article's own lines, in the article's own notation, and no retained line is substituted or re-flowed.  Retained uncounted as the source-local context the proof needs: lines 207--208.  Nothing was omitted from any retained span, so the package carries no omission record.  No retained line contains a citation, so the wrapper carries no bibliography and no bibliography entry.  No retained line contains a figure environment, an image-inclusion command or drawing code, so no image is delivered and the package declares no asset dependency.  Editorial formatting only, in addition: an AMS \texttt{amsart} preamble that restates the article's own declarations and macros used by the retained lines, namely the environments \texttt{lemma}, \texttt{theorem} on separate counters, together with the standard packages those lines need.  The retained environments print in this document as Lemma 1, Theorem 1 in order of appearance, whereas the article prints the same items with the numbering of its own full document.  The complete native source of the article as distributed in the arXiv e-print for this exact version is bundled unchanged as \texttt{sources/199961/source0.tex}.
\begin{lemma}[Proximity Principle] \label{PP} Suppose $U$ is a Kirch-open set, $n\in \mathbb N,$ and $p$ is a prime. Suppose there is some $u\in U$ such that $u \equiv n \pmod p.$ Then for every $k\in \mathbb N$ there are infinitely many $u \in U$ such that $u \equiv n \pmod{p^k}.$
\end{lemma}

\begin{theorem} \label{MainH1} Suppose $U$ and $W$ are open sets, and $V=U\cap W.$ Then every LIP function on $V$ can be written as a difference of restrictions of a LIP function on $U$ and a LIP function on $W.$
\end{theorem}

\begin{proof}
We can assume $V \neq \emptyset ,$ so all three sets are infinite. Suppose $U=\{u_1,u_2,\dots \},$ $U_n=\{u_1,u_2,\dots, u_n\}$ and similarly for $V$ and $W$. Then we have the following lemma relating splits of these finite sets.

\begin{lemma} For every $n$ for all large enough $m$ $[V_m] \in ( [V_{m+1}], [U_n], [W_n]).$
\end{lemma}

\begin{proof} For large enough $m$  $V_m$ contains $U_n\cap V$ and $W_n\cap V.$ So $U_n\setminus V_m= U_n\setminus V$ and $W_n\setminus V_m= W_n\setminus V$. Suppose $R$ is the resultant of $[U_n\setminus V]$ and $[W_n\setminus V]$. Clearly $R\neq 0$ and $R\in [U_n\setminus V], [W_n \setminus V]).$ Therefore, $R[V_m] \in ([U_n],[W_n]).$  Recall that $[V_{m+1}](x)=[V_m](x) (x-v_{m+1}).$ So the ideal $[  ( [V_{m+1}], [U_n], [W_n]) \ : \ ([V_m] )]$ contains $R$ and $x-v_{m+1}$; we want to show that it contains $1$. To do this, it is sufficient to find for every prime $p$ dividing $R$ a polynomial $f(x)$ such that $f \cdot [V_m] \in ( [V_{m+1}], [U_n], [W_n])$ and $p$ does not divide $ f(v_{m+1})$.

For every $u_i\in U_n\setminus V$ there are two possibilities: there is $v\in V$ congruent to $u_i$ modulo $p$ or there is not. We will call those $i$ that such $v$ exists, of first kind and the others of second kind (relative to $p$). Because $V$ is open in Kirch topology, by the Proximity Principle (Lemma \ref{PP}) for $i$ of the first kind one can find $v_{i,p}$ that are congruent to $u_i$ modulo $p^{ord_p(R)}.$ Moreover, we can take all $v_{i,p}$ to be distinct and not in $U_n$. Fix this collection of $\{v_{i,p}\}$. If $m$ is large enough so that $V_m$ contains all  $v_{i,p},$ we take $f(x) = \prod \limits_{i\ of\ 2nd\  kind} (x-u_i). $

Note that $p$ does not divide $f(v_{m+1})$. Note also that $f(x)\cdot [V_m](x)$ is congruent modulo $p^{ord_p(R)}\mathbb Z[x]$ to a multiple of $[U_n](x).$ Replacing $f(x)$ by $\frac{R}{p^{ord_p(R)} }f(x),$ we get that $f(v_{m+1})$ is still not divisible by $p$ and $f(x)\cdot [V_m](x) \in ( [V_{m+1}], [U_n], [W_n]),$ as promised.
\end{proof}


Due to the above lemma, we can take a strictly increasing sequence  $\{m_n\}_{n=1}^{\infty}$ such that for each $n$ for all $m\geq m_n$ we have $[V_m]\in ([V_{m+1}],[U_n], [W_n])$. Note that by repeated application of this condition, we get that $[V_{m_n}]\in ([V_{m_{n+1}}],[U_n], [W_n]).$


Suppose $f\in \textrm{LIP} (V).$ Combining together the terms of  the product-sum for $f$, we can write it as $\sum_{n=0}^{\infty} f_n(x) [V_{m_n}](x)$ for some polynomials $f_n(x).$ 
Starting with $n=1$, we rewrite $[V_{m_n}]$ as a linear combination with polynomial coefficients of  $[U_n],$ $[W_n],$ and $[V_{m_{n+1}}]$. We do this repeatedly, for $n=1,2,\dots$ to rewrite the sum for $f(x)$ as a sum of product-sums with polynomial coefficients of $[U_n]$ and $[W_n].$ This splits the sum for $f(x)$ as sum (or difference) of LIP functions on $U$ and $W.$ Note that every element of $V$ is only affected by a finite part of this infinite procedure, so the formal equality of sums is a legitimate equality of functions on $V.$
\end{proof}

\end{document}
